\documentclass[11pt]{article}

\usepackage[margin=1in]{geometry}
\usepackage{amsmath, amssymb}
\usepackage{enumitem}
\usepackage{booktabs}
\usepackage{hyperref}
\usepackage{titlesec}


\title{\textbf{STA380 Project Proposal}\\
Decoherence in a Hadamard Quantum Random Walk on a Line}
\author{Mahan Hooshmandkhayat \\
\small \texttt{mahan.hooshmandkhayat@utoronto.ca}\\
\\
\small University of Toronto, Mississauga 
}
\date{}

\begin{document}
\maketitle

\section{Project Topic}
In this project we will create an interactive R Shiny website that explores how decoherence changes the behavior of a
discrete-time quantum random walk on a one-dimensional line. We will focus on the quantum to classical transition as the decoherence strength increases, by utilizing the Hadamard coin for the baseline walk.

\section{Simulation vs.\ Dataset}
The results will be from using pure simulations. No external dataset is required.

\section{Project Details}
\subsection*{Model (Noiseless QRW)}
First, we will introduce the baseline model that we are going to be using for this project .\\

\noindent\noindent
We will implement the Hadamard QRW on a line with Hilbert space $H=H_p\otimes H_c$, where $H_p$ represents the position Hilbert space which is spanned by $\{|x\rangle\,:\,x\in \mathbb{Z}\}$, and $H_c$ represents the coin Hilbert space which is a two-dimensional qubit with the basis $\{|\uparrow\rangle,|\downarrow\rangle\}$, where $|\uparrow\rangle$ can be interpreted as moving to the ``right", and $|\downarrow\rangle$ can be interpreted as moving to the ``left".\\

\noindent Next we will need to define the following operators:
\begin{itemize}
    \item Hadamard Gate: $H=\frac{1}{\sqrt{2}}\begin{pmatrix}1&1\\1&-1\end{pmatrix}$ which determines the rotation in the coin space.
    \item Shift Operator: $S=|\downarrow\rangle\langle\downarrow|\otimes \sum_{x\in \mathbb{Z}} |x-1\rangle\langle x|+|\uparrow\rangle\langle\uparrow|\otimes \sum_{x\in \mathbb{Z}}|x+1\rangle\langle x|$ which determines the direction of the movement when it is applied on the particle.
\end{itemize}

\noindent At each step of the Hadamard walk, the wavefunction is acted on by the time evolution operator, which is given by the unitary operator, $U=S(H\otimes I_p)$ where $I_p$ represents the identity on the position space. \\

\noindent Therefore, the time evolution for a t-step walk can be given by, $|\psi_t\rangle = U^t|\psi_0\rangle$.\\


\subsection*{Decoherence (Noisy Coin)}
In the noiseless model, the particle state remains pure and evolves unitarily as $|\psi_{t+1}\rangle = U |\psi_t\rangle$. However, decoherence represents interactions with an environment, which generally drives the system into mixed states. Therefore, we have to work with the density matrix,

\[
\rho_t=|\psi_t\rangle\langle\psi_t|,\quad\text{Pure State},\qquad \rho_t=\sum_ip_i|\psi_t\rangle\langle\psi_t|,\quad\text{Mixed State}
\]

\noindent We first apply the unitary QRW step $U=S(H \otimes I_p)$, and then apply a decoherence channel to the coin. The one step evolution is,

\begin{equation}\label{eq:1}
\rho_{t+1}
= \mathcal{E}_p( U \rho_t U^\dagger)
= \sum_{k} \tilde K_k U \rho_t U^\dagger \tilde K_k^\dagger
\end{equation}
where $\{K_k\}$ are Kraus operators acting on the coin Hilbert space $H_c$ and $\tilde K_k = K_k \otimes I_p$ acts on the joint space $H_c \otimes H_p$. To ensure $\mathcal{E}_p$ is a valid quantum channel, the Kraus operators must satisfy the completeness condition $\sum_k K_k^\dagger K_k = I_c$.\\

\noindent After we have computed $\rho_t$, we can obtain the probability of observing the particle at position $x$ and at time $t$, by marginalizing over the coin states:

\begin{equation}\label{eq:2}
P_t(x) = \sum_{c\in\{\uparrow,\downarrow\}} \langle c,x \mid \rho_t \mid c,x\rangle
\end{equation}

\subsection*{Monte Carlo quantum trajectories}

By using equation \eqref{eq:1} we can evolve $\rho_t$ exactly, however storing and updating a full density matrix is very computationally expensive.\\

\noindent Instead, we will approximate the same open system dynamics by simulating many pure state trajectories. After the unitary step $|\phi\rangle=U|\psi\rangle$, we sample a Kraus operator $\tilde K_k$ with probability $p_k = \|\tilde K_k|\phi\rangle\|^2$ and update the state by $|\psi\rangle \leftarrow \frac{\tilde K_k|\phi\rangle}{\sqrt{p_k}}$. \\

\noindent In a trajectory we have a pure state $|\psi\rangle$, so its position can be defined as:

\begin{equation}\label{eq:3}
P^{(j)}_t(x) = |a^{(j)}_t(x)|^2+|b^{(j)}_t(x)|^2
\end{equation}

\noindent By averaging the measurements across many trajectories we will approximate the probabilities in \eqref{eq:2}.

\begin{equation}\label{eq:4}\hat P_t(x) =\frac{1}{N}\sum^N_{j=1}P^{(j)}_t(x)
\end{equation}

\subsection*{Decoherence Channel Types}

We will include the following two channels, each controlled by a noise level $p\in[0,1]$.

\begin{itemize}
\item Dephasing:

Let $\Pi_\uparrow = |\uparrow\rangle\langle\uparrow|$ and $\Pi_\downarrow = |\downarrow\rangle\langle\downarrow|$.

A convenient Kraus representation is

\[
K_0 = \sqrt{1-p}\,I_c, \qquad
K_1 = \sqrt{p}\,\Pi_\uparrow, \qquad
K_2 = \sqrt{p}\,\Pi_\downarrow
\]

Then the coin channel is

\begin{equation}\label{eq:5}
D_p(\rho) = (1-p)\rho + p\big(\Pi_\uparrow \rho \Pi_\uparrow + \Pi_\downarrow \rho \Pi_\downarrow\big)
\end{equation}

which suppresses off diagonal coherence in the $\{|\uparrow\rangle,|\downarrow\rangle\}$ basis. When $p=0$ we get the noiseless walk, and as $p\to 1$ coherence is quickly destroyed and we get classical like behavior.

\item Depolarizing noise:

Let the Pauli matrices on the coin be

\[
X=\begin{pmatrix}0&1\\1&0\end{pmatrix},\quad Y=\begin{pmatrix}0&-i\\ i&0\end{pmatrix},\quad Z=\begin{pmatrix}1&0\\0&-1\end{pmatrix}
\]

A standard Kraus set for depolarizing noise is

\[
K_0=\sqrt{1-p}I_c,\quad K_1=\sqrt{\tfrac{p}{3}}X,\quad K_2=\sqrt{\tfrac{p}{3}}Y,\quad K_3=\sqrt{\tfrac{p}{3}}Z
\]

so that with probability $p$ the coin is randomized by a uniformly chosen Pauli error.
\end{itemize}

\subsection*{Mean and Variance}

The mean of the noisy QRW at each time $t$ is given by, $E[X_t]=\mu_t=\sum_xx\hat P_t(x)$, and the variance is given by, $Var[X_t]=\sum_x(x-\mu_t)^2\hat P_t(x)$.

\newpage
\section{Planned Outputs}
The website will provide the following outputs:
\begin{itemize}
  \item Final position distribution at step $T$ with overlays for:
  \begin{itemize}
      \item Decoherent QRW
      \item Noiseless QRW
      \item Classical random walk
  \end{itemize}
  \item Variance vs. time: plot to highlight the transition from ballistic (quantum like) to diffusive (classical like) spreading.
  \item Summary table: reporting Monte Carlo estimates of mean and variance,
  with comparisons against the noiseless and classical baselines.
\end{itemize}

\section{User Inputs (Shiny Components)}
Users will be able to modify the following:
\begin{enumerate}
  \item Setting the seed
  \item Number of steps $T$
  \item Decoherence channel type (dephasing / depolarizing)
  \item Decoherence strength $p\in[0,1]$
  \item Number of trajectories $N$ (Monte Carlo sample size)
  \item Initial coin state can be chosen from the following options:
  \begin{itemize}
      \item Symmetric: $\frac{1}{\sqrt{2}}(|\uparrow\rangle+i|\downarrow\rangle)$
      \item Right: $| \uparrow\rangle$ 
      \item Left: $| \downarrow\rangle$
  \end{itemize}
  \item Which outputs to display (distribution, variance plot, summary table)
\end{enumerate}

\newpage

\cleardoublepage
\nocite{*}
\bibliographystyle{ieeetr}
\bibliography{ref}
\end{document}

